\chapter{Installing the program}\label{ch:installing}

There are four ways to put ATE on a machine. Pick by platform, not by preference: they are not
interchangeable, and only the first two are supported in the sense that a Station should rely on
them.

\begin{center}
\footnotesize
\begin{tabular}{@{}>{\raggedright\arraybackslash}p{3.8cm}>{\raggedright\arraybackslash}p{2.4cm}>{\raggedright\arraybackslash}p{8.6cm}@{}}
\toprule
Artifact & Support & What it is \\
\midrule
\texttt{ate-<version>-}\allowbreak\texttt{windows-x64-setup.exe} & Supported & Inno Setup installer; the frozen bundle
  under \texttt{Program Files}, \texttt{ate} added to the machine PATH, \texttt{ate setup} run at
  the end \\
\texttt{ate-<version>-}\allowbreak\texttt{linux-x86\_64.tar.gz} & Best effort & The same frozen bundle plus
  \texttt{install.sh} next to it \\
\texttt{ate\_<version>}\allowbreak\texttt{\_amd64.deb} & Best effort & The frozen bundle in \texttt{/opt/ate},
  \texttt{/usr/bin/ate} linking to it, \texttt{postinst} running \texttt{ate setup} \\
\texttt{ate-<version>-py3-none-any.whl} & Best effort & ATE as a Python package; dependencies come
  from PyPI \\
\bottomrule
\end{tabular}
\end{center}

The three frozen artifacts bundle their own CPython (3.14), Qt, QtWebEngine, lupa and pyvisa-py.
They carry no Media Releases: Media Releases have their own installers and are placed in a
Test Definition directory (chapter~\ref{ch:media-releases}). ATE never ships NI-VISA; pyvisa-py
is bundled, so a LAN or USBTMC instrument works without it, while GPIB needs NI-VISA installed
separately.

\section{Windows}

The installer is per-machine and writes to the machine PATH, so it needs administrator
privileges; it refuses 32-bit Windows. A silent install, which is what the release CI does, is:

\begin{lstlisting}[language=ATEShell]
ate-1.2.3-windows-x64-setup.exe /VERYSILENT /SUPPRESSMSGBOXES /CLOSEAPPLICATIONS
\end{lstlisting}

Setup closes any ATE that is running before replacing the files, and does not restart it. An
upgrade copies the new bundle over the old one after clearing \texttt{\_internal}, so nothing from
the previous build is left behind. Config and data are not in the install directory: they are in
\texttt{\%PROGRAMDATA\%\textbackslash ATE}, and they survive every upgrade.

At the end of the install the installer runs \texttt{ate.exe setup} itself and checks its exit
code, because a Station that is not prepared is not installed in any useful sense. If that fails
you get a dialog saying so, with the words

\begin{lstlisting}[language=ATEShell]
ATE is installed, but "ate setup" exited with code 1
\end{lstlisting}

and the fix it names: run \texttt{ate setup} from an elevated command prompt. A silent install
exits non-zero in that case, so a script that ignores the exit code will not notice a Station that
was never created.

Uninstalling removes the program and nothing else. It takes its own entry back out of the machine
PATH, entry by entry, leaving every other entry and the value's type alone; the config, database,
reports and logs stay where they are.

\textbf{What is not established here:} the installer is unsigned --- there is no certificate, so
the release build ships unsigned --- and it has only ever been installed headlessly on a CI
runner. What a particular Windows machine shows for an unsigned installer, and what the wizard
looks like to a person, are not known from this repository.

\section{Linux}

The tarball is a directory and a script. Run it from the extracted tarball:

\begin{lstlisting}[language=ATEShell]
./install.sh --prefix /opt/ate --bin /usr/bin/ate
\end{lstlisting}

\texttt{--prefix} is where the bundle goes (\texttt{PREFIX} is honoured too); \texttt{--bin} is
where the \texttt{ate} command is linked, and when it is not given the script uses
\texttt{/usr/bin/ate} as root and \texttt{<prefix>/bin/ate} for anyone else. Re-running it is an
upgrade: the bundle's own contents are cleared before the new copy is written, so no file of the
previous version survives. \texttt{./install.sh --uninstall} removes the bundle and its link, and
leaves \texttt{/etc/ate} and \texttt{/var/lib/ate} alone.

Run by an ordinary user, the link lands in the prefix's own \texttt{bin} and setup prints
\texttt{running unprivileged; skipping group/ownership/ACL setup} and skips the group and ownership
pass. That is not a failure: on a Station a root run has already prepared, setup succeeds and
\texttt{install.sh} exits 0. It exits 1 only when setup cannot do its job at all --- it cannot
create \texttt{/etc/ate} or \texttt{/var/lib/ate} --- and then it says so and tells you to finish
with \texttt{sudo ate setup}.

The \texttt{.deb} is the same bundle, and its \texttt{postinst} runs \texttt{ate setup}:

\begin{lstlisting}[language=ATEShell]
apt-get install -y ./ate_1.2.3_amd64.deb
\end{lstlisting}

It depends on the system for what a bundle cannot carry: libc, GL and EGL, the six X libraries the
Qt platform plugin needs, and \texttt{fontconfig} with \texttt{fonts-dejavu-core}. The last pair
matter more than they look: QtWebEngine renders the Test Report with the system's fonts, and with
none installed it writes a valid PDF that contains no text at all.

\section{When the install leaves the Station unprepared}

Every format runs \texttt{ate setup} as its last step, and each reports a failure of that step in a
different way, so the recovery is different too.

\textbf{Windows.} The installer runs \texttt{ate.exe setup} through \texttt{Exec} with
\texttt{SW\_HIDE}, so the setup step has no window: nothing it prints ever reaches the screen, and
the only sign of a failure is the message box. Re-running it in a console is what shows the lines:

\begin{lstlisting}[language=ATEShell]
ate setup
\end{lstlisting}

\textbf{The \texttt{.deb}.} Its \texttt{postinst} exits 1 when \texttt{ate setup} fails, and dpkg
treats a failing \texttt{postinst} as an interrupted install: the files are in \texttt{/opt/ate} and
the \texttt{ate} command is on the PATH, but the data directory, the config and the database were
never created. Other packages then refuse to install until this one is configured. Fix the cause --
an unwritable \texttt{/etc/ate} or \texttt{/var/lib/ate} is the usual one -- and then finish the
install:

\begin{lstlisting}[language=ATEShell]
sudo /opt/ate/ate setup          # creates the directories, the config and the database
sudo dpkg --configure -a         # dpkg finishes the package setup it left half done
\end{lstlisting}

\texttt{apt remove} removes only the program and keeps the config and the data. \texttt{apt purge}
also removes \texttt{/etc/ate/config.json}, and never touches \texttt{/var/lib/ate}: the database,
its WAL files, reports and logs are records, and the package deliberately does not delete them.

\section{The wheel}

The wheel uses the machine's own Python and PyPI's PySide6, so it needs more system libraries than
the bundle does, and it does not bundle Qt. \texttt{pipx} gives it its own virtual environment:

\begin{lstlisting}[language=ATEShell]
pipx install dist/ate-1.2.3-py3-none-any.whl
ate setup
\end{lstlisting}

Nothing runs \texttt{ate setup} for a wheel install: you run it yourself, as root on Linux.

\cando{choose the artifact for a platform, and tell what an upgrade and an uninstall each take
away and what they leave behind.}
